% =====================================================================
%  CS375 Lab 10 - Demand Paging, Advanced Clock, TLB and MLFQ
%
%  HOW TO USE
%   1. Fill in your name / date below (\studentname ...).
%   2. Save each screenshot into report/screenshots/ using the exact file
%      name shown in its placeholder box (e.g. 01_build.png). PNG or JPG.
%   3. Graphs and summary tables are pulled in automatically from
%      ../output/ after you run:  make experiments && make plots
%   4. Build:  make report      (or:  cd report && pdflatex report.tex  twice)
%  Anything that is missing shows up as a labelled box, so the report always
%  compiles.
% =====================================================================
\documentclass[11pt]{article}

\newcommand{\studentname}{Santiago Soto}
\newcommand{\studentid}{M00121124}
\newcommand{\coursename}{CS375 -- Operating Systems}
\newcommand{\labdate}{\today}

\usepackage[margin=1in]{geometry}
\usepackage{graphicx}
\usepackage{float}
\usepackage{booktabs}
\usepackage{caption}
\usepackage{xcolor}
\usepackage{listings}
\usepackage{fancyhdr}
\usepackage[hidelinks]{hyperref}
\usepackage{tabularx}
\setlength{\headheight}{14pt}

\definecolor{codebg}{HTML}{F5F5F3}
\lstset{basicstyle=\ttfamily\small, backgroundcolor=\color{codebg}, frame=single,
        rulecolor=\color{gray!40}, breaklines=true, columns=fullflexible}

\pagestyle{fancy}
\fancyhf{}
\lhead{\coursename\ -- Lab 10}
\rhead{\studentname}
\cfoot{\thepage}

% ---------------------------------------------------------------------
% \screenshot{file-name-without-extension}{caption}{what to capture}
% Looks for screenshots/<name>.png, then .jpg; otherwise draws a placeholder.
% ---------------------------------------------------------------------
\newcommand{\shotbox}[2]{%
  \fbox{\parbox{0.92\linewidth}{\centering\vspace{1.2em}
    {\large\textbf{Screenshot placeholder}}\\[0.4em]
    Save as \texttt{report/screenshots/\detokenize{#1}.png}\\[0.4em]
    \small #2\vspace{1.2em}}}}
\newcommand{\screenshot}[3]{%
  \begin{figure}[H]\centering
  \IfFileExists{screenshots/#1.png}{%
    \includegraphics[width=\linewidth,height=0.45\textheight,keepaspectratio]{screenshots/#1.png}}{%
  \IfFileExists{screenshots/#1.jpg}{%
    \includegraphics[width=\linewidth,height=0.45\textheight,keepaspectratio]{screenshots/#1.jpg}}{%
    \shotbox{#1}{#3}}}
  \caption{#2}\label{fig:#1}
  \end{figure}}

% \graph{file in ../output/graphs}{caption}{width}
\newcommand{\graph}[3]{%
  \begin{figure}[H]\centering
  \IfFileExists{../output/graphs/#1}{%
    \includegraphics[width=#3\linewidth]{../output/graphs/#1}}{%
    \fbox{\parbox{0.9\linewidth}{\centering\vspace{1em}Graph \texttt{output/graphs/\detokenize{#1}} not found.\\
      Run \texttt{make experiments} and \texttt{make plots}.\vspace{1em}}}}
  \caption{#2}
  \end{figure}}

% \autotable{file in ../output/summary}
\newcommand{\autotable}[1]{%
  \IfFileExists{../output/summary/#1}{\input{../output/summary/#1}}{%
    \fbox{Table \texttt{output/summary/\detokenize{#1}} not found -- run \texttt{make experiments}.}}}

\title{\textbf{Lab 10: Simulating a Virtual Memory and CPU Subsystem\\
with Demand Paging and Advanced Clock Replacement}}
\author{\studentname\ (\studentid)\\\coursename}
\date{\labdate}

\begin{document}
\maketitle
\thispagestyle{fancy}

\begin{abstract}
This lab builds a C++ simulator of a demand-paged virtual memory system working together with a
CPU scheduler. Every memory reference in a trace goes through a configurable LRU TLB, a two-level
page table and an enhanced Clock replacement policy that uses reference and modified bits. A
thread-level scheduler runs the trace under FCFS, Round Robin, Priority and a three-queue MLFQ with
aging. The simulator logs every access, generates text and \texttt{matplotlib} timeline graphs, and
reports the TLB hit rate, page-fault rate, context and thread switches, and a fairness index. The
experiments show that the scheduling policy changes memory performance as much as the hardware
parameters do. Policies that interleave many threads destroy TLB and frame locality, and a TLB
larger than the set of resident frames brings no further benefit.
\end{abstract}

\tableofcontents
\newpage

% =====================================================================
\section{Building and Running}
% =====================================================================
The project builds with a single \texttt{make}. The simulator binary is \texttt{bin/vmsim}. With no
arguments it runs all four scheduling policies on \texttt{input/trace.txt}.

\begin{lstlisting}
make                      # build bin/vmsim from src/
./bin/vmsim               # all four policies, TLB = 8, 16 frames
./bin/vmsim --policy mlfq --tlb 4 --verbose   # one policy, full log on screen
./bin/vmsim --config config/sim.conf          # settings from a config file
make experiments          # every pattern / policy / TLB size / frame count
make plots                # PNG timelines + comparison graphs
\end{lstlisting}

\screenshot{01_build}{Building the simulator with \texttt{make}.}
  {Run \texttt{make clean \&\& make} and capture the compile lines and the final link of \texttt{bin/vmsim}.}

\screenshot{02_help}{Command-line options. The TLB size, frame count, page size and policy are all configurable.}
  {Run \texttt{./bin/vmsim --help} and capture the whole usage text.}

% =====================================================================
\section{Design and Implementation}
% =====================================================================
\begin{table}[H]\centering\small
\begin{tabularx}{\linewidth}{lX}
\toprule
File & Responsibility \\
\midrule
\texttt{main.cpp} & driver: parses options, runs each policy, writes logs, CSVs and the summary \\
\texttt{Process.h/cpp} & process with threads; per-thread access streams, bursts and priorities \\
\texttt{CPU.h/cpp} & executes one cycle; counts context vs.\ thread switches \\
\texttt{Scheduler.h/cpp} & FCFS, RR (q=4), Priority, MLFQ (3 queues, aging) \\
\texttt{TLB.h/cpp} & configurable-size TLB with LRU eviction and invalidation \\
\texttt{TwoLevelPageTable.h/cpp} & outer/inner page tables, on-demand second-level tables \\
\texttt{Clock.h/cpp} & enhanced second-chance Clock with R and M bits \\
\texttt{MMU.h/cpp} & translation path TLB $\rightarrow$ page walk $\rightarrow$ page fault \\
\texttt{TraceLoader}, \texttt{Config}, \texttt{Logger}, \texttt{Report} & trace parsing, options, log file, tables and timelines \\
\bottomrule
\end{tabularx}
\caption{Source files (all in \texttt{src/}).}
\end{table}

\subsection{Two-level page table and page-fault handling}
Virtual addresses are 16 bits with 256-byte pages, so every process has 256 virtual pages. The
virtual page number is split evenly into an outer and an inner index:

\begin{table}[H]\centering
\begin{tabular}{cccl}
\toprule
bits 15--12 & bits 11--8 & bits 7--0 & example \texttt{0x1A43} \\
\midrule
outer index (4) & inner index (4) & offset (8) & page 26 = outer 1, inner 10, offset \texttt{0x43}\\
\bottomrule
\end{tabular}
\end{table}

\texttt{firstLevel} maps \texttt{(pid, outer)} to the id of a second-level table, and
\texttt{secondLevel} maps \texttt{(table id, inner)} to a frame number, where $-1$ means not
present. Second-level tables are allocated only when the first page in their range faults in. In
the FCFS run of \texttt{trace.txt}, 24 second-level tables (384 entries) were allocated, compared
with 1024 entries for four flat page tables. That is the space saving that hierarchical paging is
designed for.

Every access follows the same path in \texttt{MMU::access()}:
(1) TLB lookup;
(2) on a miss, walk both levels;
(3) if the page is not present, a page fault occurs. Clock selects a frame. If the victim is dirty
it is written back, and the victim's page-table entry and TLB entry are both invalidated. Then both
levels are updated for the new page and the TLB is refilled. Invalidating the TLB entry on
eviction matters: without it, a TLB hit could return a frame that now belongs to another page.

\screenshot{04_page_table}{Per-access log: each address is split into outer/inner/offset, and the log shows TLB misses, page-table hits and page faults. Below it are the page-table statistics.}
  {Run the two \texttt{grep} commands for screenshot 04 (page-fault log lines + ``Two-level page table'' stats).}

\subsection{Configurable TLB with LRU}
The TLB size is set with \texttt{--tlb N} or with \texttt{tlb\_size} in \texttt{config/sim.conf}.
Setting it to 0 disables the TLB. Each entry is tagged with the PID and stores a \texttt{lastUsed}
timestamp. A hit refreshes the timestamp, and an insert into a full TLB replaces the entry with the
smallest timestamp, which is true LRU. With \texttt{--tlb-flush} the TLB is flushed on every
context switch, which models an untagged TLB.

\screenshot{05_tlb}{The same trace with a 2-entry TLB and with a 16-entry TLB. Only the TLB hit rate changes; the page-fault rate stays the same.}
  {Run the two TLB commands for screenshot 05 (\texttt{--tlb 2} and \texttt{--tlb 16}) and capture both summary tables.}

\subsection{Advanced Clock with modified bits}
Each frame keeps a reference bit R and a modified bit M. Loading a page sets R=1, and M=1 if the
first access is a write. Every hit sets R, and a write also sets M. When memory is full, the hand
searches as follows:
\begin{enumerate}
\item \textbf{Pass A} -- one full sweep for a frame with (R=0, M=0) without changing any bits.
\item \textbf{Pass B} -- one full sweep for (R=0, M=1), clearing R on every frame it passes.
\item Repeat. After pass B every R bit is 0, so the next round always finds a victim.
\end{enumerate}
Clean, unreferenced pages are therefore evicted first, and a dirty page is written back only when
no clean candidate exists. The trace \texttt{input/clock\_demo.txt} demonstrates this with 3 frames.
When page 3 arrives, the hand reaches dirty page 0 first, but clean page 1 is evicted instead. In
the full \texttt{trace.txt} FCFS run, 47 of the 99 evictions were clean and needed no disk write.

\screenshot{06_clock}{Clock demo: clean pages are evicted before dirty ones. The final frame table shows the R and M bits.}
  {Run the clock demo command for screenshot 06 and capture the log lines, the Clock statistics and the frame table.}

\subsection{Thread simulation and scheduling (FCFS, RR, Priority, MLFQ)}
Each trace line carries a PID and a thread ID. A \texttt{Process} keeps a separate access stream,
CPU burst (number of accesses) and priority for each thread, so the schedulable unit is a thread
\texttt{(pid, tid)}. One tick is one CPU cycle, and each cycle performs one memory access. The CPU
distinguishes two kinds of switch:
\begin{itemize}
\item \textbf{context switch} -- the next thread belongs to a \emph{different process}, so the
      address space changes;
\item \textbf{thread switch} -- the next thread belongs to the \emph{same process}, so the address
      space and the TLB contents stay valid.
\end{itemize}
\textbf{FCFS} runs threads to completion in arrival order. \textbf{RR} uses a quantum of 4.
\textbf{Priority} always runs the ready thread with the lowest priority number and shares the CPU
round robin among equal priorities. Process and thread priorities come from \texttt{PRIORITY} and
\texttt{THREAD} lines in the trace. \textbf{MLFQ} has three queues with quanta 4, 8 and 16. Every
thread starts in Q0 and is demoted one level each time it uses its whole quantum. A thread that has
waited 30 ticks in Q1 or Q2 is promoted one level (aging), which prevents starvation.

\screenshot{07_mlfq}{MLFQ log: demotions after a full quantum, promotions by aging, and the per-thread table with demotion and promotion counts.}
  {Run the two MLFQ commands for screenshot 07.}

\screenshot{08_threads}{Thread-level scheduling under RR: thread switches (same address space) and context switches (new process), plus per-thread burst, wait and turnaround times.}
  {Run the two thread commands for screenshot 08.}

% =====================================================================
\section{Modified Trace and Access Patterns}
% =====================================================================
\texttt{scripts/gen\_traces.py} uses a fixed seed, so every result in this report can be reproduced
exactly. The modified \texttt{input/trace.txt} keeps the 7 original sample lines and adds four
sections: P1 \emph{sequential} (a strided array walk, 8 accesses per page), P2 \emph{random}
(uniform over 64\,KiB), P3 \emph{localized} (90\% of accesses in a 3-page hot set that moves
halfway through), and P4 \emph{mixed} (one localized thread and one sequential thread). The three
pattern traces \texttt{trace\_sequential/random/localized.txt} use the \emph{same} 8 threads, burst
lengths and priorities, so the access pattern is the only variable. About 30\% of all accesses are
writes.

\screenshot{03_trace}{The modified \texttt{input/trace.txt}: header, priorities, the original sample and the start of the sequential section.}
  {Run \texttt{head -n 45 input/trace.txt}.}

% =====================================================================
\section{Results}
% =====================================================================
Default configuration: TLB = 8 entries (LRU, tagged), 16 frames, 256-byte pages, quantum 4, aging
30. All 323 accesses of \texttt{trace.txt} take 323 ticks under every policy because there is no
idle time, so the policies differ only in \emph{order}. Even so, the order changes memory
performance considerably.

\subsection{Summary table of performance metrics}
\begin{table}[H]\centering\small
\autotable{trace_table.tex}
\caption{\texttt{input/trace.txt} (mixed patterns), generated automatically by the simulator. Fairness is Jain's index of $\textit{burst}/\textit{turnaround}$ over all threads (1 = perfectly fair).}
\end{table}

\begin{table}[H]\centering\small
\textbf{Sequential}\par\smallskip
\autotable{trace_sequential_table.tex}\par\bigskip
\textbf{Random}\par\smallskip
\autotable{trace_random_table.tex}\par\bigskip
\textbf{Localized}\par\smallskip
\autotable{trace_localized_table.tex}
\caption{The same 8 threads with a single access pattern each.}
\end{table}

\screenshot{10_summary}{Terminal output of the performance summary for \texttt{trace.txt}.}
  {Run \texttt{./bin/vmsim --quiet} and capture the summary table.}

\subsection{Timeline graphs}
Each tick is coloured by the outcome of that access. Grey bars show the time a thread spends ready
but waiting, solid vertical lines mark context switches and dotted lines mark thread switches.

\graph{timeline_trace_fcfs.png}{FCFS timeline: long uninterrupted bursts, so pages are faulted in once and then hit in the TLB.}{1.0}
\graph{timeline_trace_rr.png}{Round Robin timeline: a context switch every 4 ticks, with TLB misses and faults at the start of almost every quantum.}{1.0}
\graph{timeline_trace_priority.png}{Priority timeline: the priority-1 threads finish first and the priority-3 threads (P3, P1.T2) wait until the end.}{1.0}
\graph{timeline_trace_mlfq.png}{MLFQ timeline: quanta grow from 4 to 16 as threads are demoted, and aging promotions keep the waiting threads moving.}{1.0}

\screenshot{09_ascii_timeline}{The text-based (ASCII) timeline printed by the simulator: one row per thread, plus a row for context and thread switches.}
  {Run the ASCII timeline command for screenshot 09.}

\subsection{Access patterns, TLB size and frame count}
\graph{policy_comparison.png}{Policy comparison on \texttt{trace.txt}.}{0.95}
\graph{pattern_comparison.png}{TLB hit rate and page-fault rate for each access pattern and policy.}{0.95}
\graph{tlb_sweep.png}{TLB hit rate as the TLB grows from 1 to 32 entries (16 frames).}{1.0}
\graph{frames_sweep.png}{Page-fault rate as the number of physical frames grows (\texttt{trace.txt}).}{0.65}

\screenshot{11_patterns}{\texttt{make experiments}: summary tables for the sequential, random and localized traces.}
  {Run \texttt{make experiments} and capture the first summary tables it prints.}

% =====================================================================
\section{Analysis}
% =====================================================================
\subsection{Scheduling policy vs.\ memory performance}
\paragraph{TLB hit rate.} FCFS has the best TLB hit rate on the mixed trace (63.8\%), and RR has the
worst (39.0\%). Priority (62.2\%) is close to FCFS, and MLFQ falls in between (46.4\%). The
explanation is the \emph{reuse distance} between two accesses to the same page. Under FCFS a thread
uses its pages in one uninterrupted burst, so its working set stays in the TLB. Under RR, all eight
threads take turns every 4 ticks. By the time a thread runs again, seven other threads have each
inserted one to three translations, and an 8-entry LRU TLB has evicted most of its entries. The
RR timeline shows a miss or a fault at the start of almost every quantum. Priority behaves like
FCFS because, within each priority level, only two or three threads compete. MLFQ starts as RR
(quantum 4), but demoted threads get quanta of 8 and 16 and reuse their pages for longer.

\paragraph{Page-fault rate.} The same effect appears one level lower. With 16 frames, FCFS and
Priority fault on about 35\% of accesses, while RR and MLFQ fault on about 48\%. Interleaving all
eight threads means their combined working sets (about 24 hot pages for the localized threads
alone) compete for 16 frames at the same time, so Clock evicts pages that will be needed again a
few quanta later. Faults also cause disk writes: RR needed 60 dirty write-backs and FCFS needed 52.
The frame sweep shows that the RR/MLFQ disadvantage shrinks as memory grows, from 64.7\% vs.\ 37.8\%
faults with 4 frames to 38.7\% vs.\ 34.7\% with 32. With enough memory, ordering matters much less,
because what remains are compulsory faults.

\paragraph{Context-switch overhead.} FCFS performs only 5 context switches and 2 thread switches.
RR performs 67 context switches, more than 13 times as many, with one at almost every quantum
boundary. Priority (42) and MLFQ (41) are in between. MLFQ saves switches through its longer
lower-level quanta, and Priority saves them because only a few threads share a level. The
difference between a \emph{thread} switch and a \emph{context} switch matters here. A thread switch
keeps the address space, so a tagged TLB and resident pages stay useful. A context switch changes
it. The tagged-versus-flushed experiment measures the cost directly: flushing the TLB on each
context switch leaves FCFS unchanged (63.8\%, only 5 switches), but lowers RR from 39.0\% to 34.7\%
and Priority from 62.2\% to 44.0\%. PID-tagged TLBs exist precisely to avoid this cost.

\paragraph{Thread scheduling fairness.} The fairness ranking is the reverse of the memory ranking.
Jain's index is 0.576 for FCFS, 0.715 for Priority, and 0.943 and 0.947 for RR and MLFQ. Under FCFS,
P4.T2 has a burst of only 12 ticks but finishes at tick 323 after waiting 311 ticks, because it
arrived last. RR and MLFQ give every thread a first response within 28 ticks (average response
time 14.0 versus 141.2 for FCFS). Priority starves the priority-3 threads: P3.T1 and P3.T2 wait 273
and 274 ticks and do not get the CPU until ticks 220 and 224. MLFQ combines a fast response for
everyone with fewer switches than RR. Its 39 demotions and 29 aging promotions keep long CPU-bound
threads such as P4.T1 (burst 80) from monopolising the CPU without starving anyone. The aging
threshold is a real trade-off. At 15 ticks, MLFQ makes 57 context switches and behaves almost like
RR. At 60 ticks or more it makes only 26 switches and its TLB hit rate rises to 51.7--53.6\%, at a
small cost in fairness (0.932--0.938).

\paragraph{Overall.} No single policy wins on every metric. FCFS and Priority give the best TLB and
page-fault rates, at a severe cost in response time and fairness. RR is the fairest but is the
worst for the memory system. MLFQ is the best compromise: its fairness is as high as RR's, it makes
39\% fewer context switches, and it has a better TLB hit rate. These measurements reproduce the
classic tension between locality and interactivity.

\subsection{Impact of memory access patterns}
The access pattern matters more than the scheduler. With identical threads and the default
configuration:
\begin{itemize}
\item \textbf{Sequential} access gives a TLB hit rate of 87.3\% and 12.7\% faults under FCFS, RR and
      Priority. Each 256-byte page is touched 8 times (stride 32), so 7 of every 8 accesses hit and
      only the first access to each new page faults. These are compulsory faults that no
      replacement policy can avoid. Because each thread uses only one page at a time, eight threads
      need only eight translations, which fit exactly in the 8-entry TLB, so even RR loses nothing.
      MLFQ is slightly lower (82.0\%) because its 16-tick quanta cross two pages per turn and push
      out other threads' entries.
\item \textbf{Random} access defeats both caches: 95--98\% of accesses fault and the TLB hit rate is
      below 3\%. With 256 possible pages per process and 16 frames, almost no page is resident when
      it is touched again. No policy and no reasonable TLB size helps; only a much larger memory
      would.
\item \textbf{Localized} access is the most sensitive to scheduling. Under FCFS it is almost as good
      as sequential (81.0\% hits, 19.0\% faults), because each thread's 3-page hot set fits easily.
      Under RR the hit rate falls to 45.9\% and faults double to 37.7\%, because eight hot sets of
      3 pages (24 pages) do not fit in 16 frames or 8 TLB entries at the same time. Locality only
      pays off if the scheduler lets a thread run long enough to use it.
\end{itemize}

\subsection{Impact of TLB size}
The TLB-sweep graph shows three regimes:
\begin{enumerate}
\item \textbf{TLB smaller than the working set}: each doubling helps a lot. For localized/FCFS the
      hit rate goes 29.4\% $\rightarrow$ 55.7\% $\rightarrow$ 80.4\% at 1, 2 and 4 entries, because
      a 3-page hot set needs about 4 entries.
\item \textbf{TLB covers the working set}: the curve flattens. For FCFS, going from 4 to 32 entries
      gains less than 2 percentage points.
\item \textbf{TLB larger than physical memory}: no gain at all. Every curve is identical at 16 and
      32 entries, because a TLB can only map pages that are resident, and there are only 16 frames.
      Every miss beyond that point is a page fault. TLB hit rate + page-fault rate $\le 100\%$, so
      RR on localized access can never exceed $100 - 37.7 = 62.3\%$, and it reaches exactly 62.3\%.
\end{enumerate}
Interleaving policies need a larger TLB to reach the same hit rate. RR and MLFQ on the mixed trace
keep improving up to 16 entries (39.0\% $\rightarrow$ 51.7\%), while FCFS saturates at 4 entries.
Random access never exceeds 5.1\% at any size.

\subsection{Conclusion}
The simulator demonstrates demand paging with a two-level table, an LRU TLB and an enhanced Clock
that avoids disk writes by preferring clean victims. Memory performance is a property of the whole
system, not only of the MMU. The scheduling policy determines how much of each thread's locality
survives. Long, uninterrupted bursts (FCFS, Priority, or MLFQ's lower queues) preserve TLB and frame
contents. Frequent switching among many threads (RR) increases misses, faults and write-backs. A
well-tuned MLFQ with aging gets most of RR's fairness at a much lower memory cost.

\appendix
\section{Generated Files}
\screenshot{12_plotter}{\texttt{python3 scripts/plotter.py}: the ASCII graphs and the list of generated PNG files.}
  {Run \texttt{make plots | tail -n 24} and capture the pattern bars and the ``PNG graphs:'' list.}
\screenshot{13_log_files}{Auto-generated log files and the start of a log.}
  {Run \texttt{ls -lh output/logs} and \texttt{head -n 30 output/logs/trace\_mlfq.log}.}

\begin{table}[H]\centering\small
\begin{tabularx}{\linewidth}{lX}
\toprule
Path & Contents \\
\midrule
\texttt{output/logs/<trace>\_<policy>.log} & full per-tick log, statistics, frame table, ASCII timeline \\
\texttt{output/events/<trace>\_<policy>.csv} & one row per tick (input to \texttt{plotter.py}) \\
\texttt{output/summary/<trace>.csv} & summary metrics, one row per policy \\
\texttt{output/summary/*\_sweep.csv} & TLB-size and frame-count sweeps, tagged vs.\ flushed TLB \\
\texttt{output/graphs/*.png} & timeline and comparison graphs (plus \texttt{ascii\_graphs.txt}) \\
\bottomrule
\end{tabularx}
\end{table}

\end{document}
